% Author: Benjamin Stanley Frohman
% Affiliation: Independent Researcher, Texas, United States
%
\documentclass[11pt,a4paper]{article}

\usepackage[margin=1in]{geometry}
\usepackage{amsmath,amssymb,amsthm}
\usepackage{xurl}
\usepackage{hyperref}
\emergencystretch=3em

\newtheorem{theorem}{Theorem}[section]
\newtheorem{proposition}[theorem]{Proposition}
\newtheorem{remark}[theorem]{Remark}

\title{Strain-Tensor Eigenvalue Decomposition\\and Intermediate Alignment in 3D Navier--Stokes}
\author{Benjamin Stanley Frohman\thanks{\textcopyright\ 2026 Benjamin Stanley Frohman. Distributed under the Apache License 2.0.}\\\small Independent Researcher, Texas, United States\\\small ORCID \href{https://orcid.org/0009-0006-7068-3718}{0009-0006-7068-3718}\\\small \href{mailto:frohmanbenjamin@gmail.com}{frohmanbenjamin@gmail.com}}
\date{10 November 2026}

\begin{document}
\maketitle

\begin{abstract}
We record the elementary algebraic decomposition of the symmetric strain-rate tensor of a divergence-free velocity field and its consequences for the local vortex-stretching term. Because the trace of the strain vanishes, the intermediate eigenvalue controls the leading contribution once the vorticity direction is aligned with the intermediate eigenvector. The identities are pointwise and unconditional, and the trace identity, the ordering \(\lambda_1\le\lambda_2\le\lambda_3\), and the intermediate-eigenvalue representation are checked in Lean~4 (Section~\ref{sec:lean}); the geometric and analytic discussion is not formalized. The identities do not by themselves produce an a priori bound on the \(L^\infty\) norm of the strain and therefore leave the supercritical gap (Hypothesis~B) open.
\end{abstract}

\section{Eigenvalue structure of the strain tensor}

Let \(u\) be a smooth divergence-free vector field on \(\mathbb{T}^3\) and let
\[
S=\frac12\bigl(\nabla u+(\nabla u)^T\bigr)
\]
be the symmetric strain-rate tensor. Then \(\operatorname{tr}S=\nabla\cdot u=0\), so the three real eigenvalues satisfy (cf.\ Betchov~\cite{Betchov1956})
\begin{equation}\label{eq:trace}
\lambda_1\le\lambda_2\le\lambda_3,\qquad\lambda_1+\lambda_2+\lambda_3=0.
\end{equation}
At each point where \(S\) is continuous there exists an orthonormal eigenbasis \(\{e_1,e_2,e_3\}\). Writing the unit vorticity direction
\[
\boldsymbol{\xi}=\frac{\omega}{|\omega|}=\xi_1 e_1+\xi_2 e_2+\xi_3 e_3
\]
(with \(\xi_1^2+\xi_2^2+\xi_3^2=1\)) one obtains the quadratic form
\begin{equation}\label{eq:quadratic}
\boldsymbol{\xi}\cdot S\boldsymbol{\xi}=\lambda_1\xi_1^2+\lambda_2\xi_2^2+\lambda_3\xi_3^2.
\end{equation}

\begin{proposition}[Intermediate-eigenvalue representation]
\[
\boldsymbol{\xi}\cdot S\boldsymbol{\xi}
=\lambda_2+(\lambda_1-\lambda_2)\xi_1^2+(\lambda_3-\lambda_2)\xi_3^2.
\]
\end{proposition}

\begin{proof}
Substitute \(\xi_2^2=1-\xi_1^2-\xi_3^2\) into~\eqref{eq:quadratic} and rearrange.
\end{proof}

\section{Geometric interpretation}

Because \(\lambda_1\le 0\le\lambda_3\), the deviation terms \((\lambda_1-\lambda_2)\xi_1^2\) and \((\lambda_3-\lambda_2)\xi_3^2\) measure the departure of \(\boldsymbol{\xi}\) from the intermediate axis \(e_2\). Direct numerical simulations (Ashurst, Kerstein, Kerr and Gibson~\cite{AKKG1987}) show that intense stretching in turbulent flows correlates with preferential alignment of \(\boldsymbol{\xi}\) with \(e_2\) rather than with the most extensional eigenvector \(e_3\). Regularity criteria involving only the positive part \(\lambda_2^+\) of the intermediate eigenvalue are due to Neustupa and Penel~\cite{NeustupaPenel2001,NeustupaPenel2003,NeustupaPenel2005} and Miller~\cite{Miller2020,Miller2020Addendum}. The algebraic identity above makes this statement precise: the leading contribution is exactly \(\lambda_2\), while the remainder is controlled by the angular dispersion of \(\boldsymbol{\xi}\).

\section{Relation to the alignment kernel}

Under a Hölder condition
\[
|\boldsymbol{\xi}(x)-\boldsymbol{\xi}(y)|\le K|x-y|^\alpha,\qquad\alpha\ge\frac12,
\]
on high-vorticity sets (Constantin and Fefferman~\cite{ConstantinFefferman1993} treat the Lipschitz case \(\alpha=1\); the relaxation to \(\alpha\ge\frac12\) is due to Beir\~ao da Veiga and Berselli~\cite{BeiraoBerselli2002}), the Constantin--Fefferman singular integral that represents \(\boldsymbol{\xi}\cdot S\boldsymbol{\xi}\) desingularizes. Heuristically, the eigenvalue decomposition isolates the same phenomenon; we do \emph{not} prove here that the integral of the deviation terms is controlled by the angular dispersion functional
\[
\mathcal{D}_{\boldsymbol{\xi}}(t)=\iint_{\Omega_M\times\Omega_M}|\omega(x)|^2|\omega(y)|^2|\boldsymbol{\xi}(x)\times\boldsymbol{\xi}(y)|^2\,K_{\mathrm{BS}}(x-y)\,dx\,dy,
\]
and this comparison is stated only as motivation.

\section{Conditional status}

The identities of this note are pointwise and algebraic. They do not imply integrability of \(\|S\|_{L^\infty}\) in time. In particular, they do not prove Hypothesis~B
\[
\int_0^T\|S(t)\|_{L^\infty}\,dt<\infty
\]
for arbitrary large smooth data; a time-integrable bound on \(\|S\|_{L^\infty}\) would require further hypotheses that remain open. The note makes no claim about the Clay Millennium problem. The supercritical gap between the energy inequality and the Beale--Kato--Majda criterion remains open.


\section{Lean formalization: exact scope}\label{sec:lean}

The repository \url{https://github.com/BenFrohman/frohmanian-strain-eigenvalue} contains \texttt{StrainEigenvalue.lean} (Lean~4.30.0-rc1, Mathlib \texttt{v4.30.0-rc1}, commit \texttt{0c154d6}). A \texttt{StrainTensor} is a real symmetric \(3\times3\) matrix with zero trace. Lean indices \(0,1,2\) play the roles of \(1,2,3\) here. The public eigenvalue family is Mathlib's antitone family \texttt{IsHermitian.eigenvalues\textsubscript{0}}, reversed, so the eigenvalues are increasing: Lean proves \(\lambda_1\le\lambda_2\le\lambda_3\) (\texttt{eigenvalues\_ordered}) and, with the trace identity, \(\lambda_1\le0\le\lambda_3\) (\texttt{eigenvalues\_min\_nonpos\_max\_nonneg}). Index \(2\) in the paper's labelling (Lean index \(1\)) is therefore the intermediate eigenvalue. Eigenvectors and eigen-coordinates use the same order. The kernel checks the following, using only the axioms \texttt{propext}, \texttt{Classical.choice}, \texttt{Quot.sound}:
\begin{itemize}
\item \texttt{eigenvalues\_ordered}: \(\lambda_1\le\lambda_2\le\lambda_3\).
\item \texttt{eigenvalues\_sum\_zero}: \(\lambda_1+\lambda_2+\lambda_3=0\).
\item \texttt{eigenvalues\_min\_nonpos\_max\_nonneg}: \(\lambda_1\le0\le\lambda_3\).
\item \texttt{quadratic\_form\_eq\_sum}: with \(U\) the orthonormal eigenvector matrix in this increasing order and \(\eta=U^{T}\xi\), \(\xi\cdot S\xi=\sum_i\lambda_i\eta_i^2\) for every \(\xi\in\mathbb{R}^3\) (equation~\eqref{eq:quadratic}).
\item \texttt{sum\_sq\_eigenCoords}: \(\sum_i\eta_i^2=\sum_i\xi_i^2\).
\item \texttt{quadratic\_form\_intermediate}: if \(\sum_i\xi_i^2=1\), then \(\xi\cdot S\xi=\lambda_2+(\lambda_1-\lambda_2)\eta_1^2+(\lambda_3-\lambda_2)\eta_3^2\), with \(\lambda_2\) the intermediate eigenvalue. This is Proposition~1.1.
\item Non-vacuity: \(S=\operatorname{diag}(1,-1,0)\) is a \texttt{StrainTensor} whose increasing eigenvalues are \((-1,0,1)\) (\texttt{exampleStrain\_eigenvalues\_eq}); they are not all zero, and for \(\xi=(1,0,0)\) both sides of Proposition~1.1 equal \(1\).
\end{itemize}
The theorem \texttt{intermediate\_representation} proves the same identity for \texttt{stretching}, which is \(\sum_i\lambda_i\xi_i^2\) with \(\xi\) already in ordered eigen-coordinates. An earlier skeleton gave \texttt{stretching} the hypothesis \(\|\xi\|=1\). On \(\mathrm{Fin}\,3\to\mathbb{R}\), Mathlib's \(\|\cdot\|\) is the sup norm, so that does not follow from \(\sum_i\xi_i^2=1\): \(\xi=(3/5,4/5,0)\) has \(\sum\xi_i^2=1\) but \(\|\xi\|_\infty=4/5\). Lean checks this counterexample (\texttt{sum\_sq\_eq\_one\_not\_imp\_norm\_eq\_one}). The hypothesis is now \(\sum_i\xi_i^2=1\). The development has no \texttt{sorry}.

Nothing about vorticity, Navier--Stokes dynamics, H\"older alignment, the Constantin--Fefferman kernel, \(\mathcal{D}_{\boldsymbol\xi}\), or Hypothesis~B is formalized. A kernel check establishes the algebraic statements above and nothing more.

\begin{thebibliography}{99}
\bibitem{ConstantinFefferman1993} P. Constantin and C. Fefferman, Direction of vorticity and the problem of global regularity for the Navier--Stokes equations, Indiana Univ. Math. J. \textbf{42} (1993), 775--789. \href{https://doi.org/10.1512/iumj.1993.42.42034}{doi:10.1512/iumj.1993.42.42034}.
\bibitem{Betchov1956} R. Betchov, An inequality concerning the production of vorticity in isotropic turbulence, J. Fluid Mech. \textbf{1} (1956), no.~5, 497--504. \href{https://doi.org/10.1017/S0022112056000317}{doi:10.1017/S0022112056000317}.
\bibitem{AKKG1987} W. T. Ashurst, A. R. Kerstein, R. M. Kerr and C. H. Gibson, Alignment of vorticity and scalar gradient with strain rate in simulated Navier--Stokes turbulence, Phys. Fluids \textbf{30} (1987), no.~8, 2343--2353. \href{https://doi.org/10.1063/1.866513}{doi:10.1063/1.866513}.
\bibitem{BeiraoBerselli2002} H. Beir\~ao da Veiga and L. C. Berselli, On the regularizing effect of the vorticity direction in incompressible viscous flows, Differential Integral Equations \textbf{15} (2002), no.~3, 345--356. \href{https://doi.org/10.57262/die/1356060864}{doi:10.57262/die/1356060864}.
\bibitem{NeustupaPenel2001} J. Neustupa and P. Penel, Anisotropic and geometric criteria for interior regularity of weak solutions to the 3D Navier--Stokes equations, in \emph{Mathematical Fluid Mechanics: Recent Results and Open Questions}, Birkh\"auser, Basel, 2001, pp.~237--265. \href{https://doi.org/10.1007/978-3-0348-8243-9_10}{doi:10.1007/978-3-0348-8243-9\_10}.
\bibitem{NeustupaPenel2003} J. Neustupa and P. Penel, The role of eigenvalues and eigenvectors of the symmetrized gradient of velocity in the theory of the Navier--Stokes equations, C. R. Math. Acad. Sci. Paris \textbf{336} (2003), no.~10, 805--810. \href{https://doi.org/10.1016/S1631-073X(03)00174-2}{doi:10.1016/S1631-073X(03)00174-2}.
\bibitem{NeustupaPenel2005} J. Neustupa and P. Penel, Regularity of a weak solution to the Navier--Stokes equation in dependence on eigenvalues and eigenvectors of the rate of deformation tensor, in \emph{Trends in Partial Differential Equations of Mathematical Physics}, Progr. Nonlinear Differential Equations Appl. \textbf{61}, Birkh\"auser, Basel, 2005, pp.~197--212. \href{https://doi.org/10.1007/3-7643-7317-2_15}{doi:10.1007/3-7643-7317-2\_15}.
\bibitem{Miller2020} E. Miller, A regularity criterion for the Navier--Stokes equation involving only the middle eigenvalue of the strain tensor, Arch. Ration. Mech. Anal. \textbf{235} (2020), no.~1, 99--139. \href{https://doi.org/10.1007/s00205-019-01419-z}{doi:10.1007/s00205-019-01419-z}.
\bibitem{Miller2020Addendum} E. Miller, Addendum to: A regularity criterion for the Navier--Stokes equation involving only the middle eigenvalue of the strain tensor, Arch. Ration. Mech. Anal. \textbf{237} (2020), no.~3, 1173--1175. \href{https://doi.org/10.1007/s00205-020-01527-1}{doi:10.1007/s00205-020-01527-1}.
\end{thebibliography}

\end{document}
